% Page budget: 3 pages
\chapter{מבוא}

\section{מגמות בתחום}
שוק כלי יצירת הווידאו מבוססי בינה מלאכותית (\textenglish{AI}) גדל בקצב דו־ספרתי~\cite{fbi2026aivideo}, והיצירה האוטומטית של וידאו הופכת מניסוי טכנולוגי למוצר מסחרי [TO-VERIFY].לפי חברת הניתוח \textenglish{Fortune Business Insights}, היקף שוק מחוללי הווידאו מבוססי \textenglish{AI} הוערך ב־\textenglish{716.8} מיליון דולר ב־\textenglish{2025}, צפוי להגיע ל־\textenglish{847} מיליון דולר ב־\textenglish{2026}, ולצמוח ל־\textenglish{3,350} מיליון דולר ב־\textenglish{2034}, בקצב צמיחה שנתי מצטבר (\textenglish{CAGR}) של \textenglish{18.8}\% בתקופה \textenglish{2026--2034}~\cite{fbi2026aivideo}. צפון אמריקה החזיקה בכ־\textenglish{41}\% מהשוק ב־\textenglish{2025}~\cite{fbi2026aivideo}. נתון זה מעיד על ריכוז גיאוגרפי משמעותי של השוק ביבשת אחת.

הערכות השוק אינן אחידות. דוח ישן וצר יותר של \textenglish{Markets\-and\-Markets} הגדיר את שוק הטקסט־לווידאו (\textenglish{Text-to-Video}) כ־\textenglish{0.1} מיליארד דולר ב־\textenglish{2022} ו־\textenglish{0.9} מיליארד דולר ב־\textenglish{2027}, עם \textenglish{CAGR} של \textenglish{37.1}\%~\cite{mnm2023t2v}. ההפרש בין הדוחות נובע מהגדרות שונות של השוק, ולכן הדוח הנוכחי מתייחס לנתונים אלה כאל טווח ולא כאל מספר יחיד. פרק~4 מפרט את ההשוואה.

במקביל גדלה כלכלת היוצרים (\textenglish{Creator Economy}). לפי הערכת \textenglish{Goldman Sachs Research}, היקפה הפוטנציאלי עשוי לעלות מכ־\textenglish{250} מיליארד דולר ב־\textenglish{2023} לכ־\textenglish{480} מיליארד דולר ב־\textenglish{2027}~\cite{goldman2023creator}. אולם רק כ־\textenglish{4}\% מהיוצרים הם מקצועיים, כלומר מרוויחים יותר מ־\textenglish{100} אלף דולר בשנה~\cite{goldman2023creator}. הנתונים מרמזים על קהל לא מבוטל של יוצרים עצמאיים. קהל זה מחפש כלים שיאפשרו להם להפיק תוכן ברמה גבוהה בלי צוות הפקה.

\section{הקשר יישומי}
התחום היישומי של הדוח הוא סדרות וידאו קצרות סביב דמויות קבועות, המופצות ביוטיוב ובטיקטוק. בסדרה כזו הצופה חוזר לפרק הבא משום שהוא מכיר את הדמויות ואת העולם, ולכן הערך של הסדרה תלוי ברציפות. הרציפות כוללת שני רבדים: עקביות חזותית, כלומר הדמות נראית זהה משוט לשוט, ורציפות עלילתית, כלומר אירועי עבר, יחסים בין דמויות ומצבן נשמרים מפרק לפרק.

קהל היעד העיקרי הוא יוצרי תוכן דוברי עברית ללא רקע טכני. הרחבה עתידית אפשרית היא משרדי פרסום, שבהם דמויות מותג וקמפיינים מתמשכים דורשים את אותה רציפות. היקף קהל היעד בישראל הוא [TO-VERIFY] יוצרים.

כיום מי שמבקש להפיק סדרה כזו משתמש בכמה כלים נפרדים. קיימים כלים מסחריים כגון \textenglish{LTX Studio}~\cite{ltxstudio}, \textenglish{OpenArt}~\cite{openart} ו־\textenglish{LongStories.ai}~\cite{longstories}, וכן תהליך ידני המשלב כלי עריכה כמו \textenglish{CapCut}~\cite{capcut}. המוצרים האלה נבדקים לעומק בפרק~4.

\section{מוטיבציה}
הפתרונות הקיימים אינם מספקים יוצר שמבקש להפעיל סדרה ארוכה בעברית, מארבע סיבות. ראשית, התהליך מפוצל: היוצר מחבר בעצמו מודל תמונה, מודל וידאו, קול ועריכה, והמידע על הדמויות אינו נשמר במקום אחד. שנית, ניהול הרציפות העלילתית מוטל על היוצר. לפי בדיקת הצוות, הכלים שומרים על מראה הדמות ועל עולמה, ואחד מהם, \textenglish{LongStories.ai}, שומר גם על הדמויות, הקולות והעולם בין פרקים~\cite{longstories}, אך אינם מתעדים את אירועי העלילה ואינם מתריעים כשפרק חדש סותר פרק קודם [TO-VERIFY]. משמעות הדבר היא שהיוצר נדרש לזכור בעצמו מה אירע בכל פרק קודם ולבדוק ידנית שהפרק החדש אינו סותר אותו. שלישית, תמיכה בעברית בקלט, בקול, בסנכרון שפתיים ובכתוביות מימין לשמאל (\textenglish{RTL}) אינה מתועדת במקור: [TO-VERIFY]. רביעית, זמן ההפקה של דקת וידאו בתהליך הידני הוא [TO-VERIFY] שעות, ועלות הדקה היא [TO-VERIFY] דולר.

המחקר האקדמי מתמקד בעיקר בעקביות חזותית של דמויות ושל סצנות. עבודות כמו \textenglish{MovieDreamer}~\cite{zhao2025moviedreamer} ו־\textenglish{StoryDiffusion}~\cite{zhou2024storydiffusion} שומרות על זהות הדמות לאורך רצף ארוך, ו־\textenglish{StoryAgent}~\cite{hu2024storyagent} מציגה צוות סוכנים להפקה. סקירה מלאה מופיעה בפרק~3. הפער שמתגלה הוא שאף אחד מהם אינו מנהל זיכרון של העלילה עצמה ואינו מזהה סתירות עלילתיות בין פרקים, ואין מערכת קצה־לקצה הנגישה ליוצרים בעברית בלי תשלום.

\begin{directornote}
הפרויקט נולד מהפער הזה: מערכת של סוכנים בינה מלאכותית הפועלת כ״מנהל סדרה״ (\textenglish{Showrunner}), המנהלת ״ספר סדרה״ שמתעדכן אחרי כל פרק, מתריעה על סתירות ומקבלת קלט בעברית. הצלחת הפרויקט תימדד ביכולת לשמור על הדמות ועל העלילה לאורך פרקים, ולא רק בבניית עוד מחולל וידאו.
\end{directornote}

\section{מבנה הדוח}
הדוח בנוי כמשפך, מהרחב אל המסוים: הוא נפתח במגמות השוק, ממשיך להקשר היישומי ולמוטיבציה, וממשיך משם אל הבעיה, הספרות, השוק, התכן ותוכנית העבודה.פרק~2 מגדיר את הבעיה בצורה כמותית ומציג יעדי תכן ומפרט דרישות. פרק~3 סוקר את הספרות באופן תמטי, מנתח מאמר דגל ומנסח את פער הידע. פרק~4 מעריך את היקף השוק ומשווה בין מוצרים מתחרים. פרק~5 משווה חלופות תכן במטריצת החלטה משוקללת, ומציג את הארכיטקטורה המוצעת ואת החדשנות. פרק~6 מציג את תוכנית העבודה לסמסטר ב׳, הכוללת חלוקת משימות ולוח זמנים של ארבעה־עשר שבועות. פרק~7 מנתח סיכונים ופעולות מנע, ופרק~8 מסכם את הממצאים ואת שלוש הפעולות המיידיות לאחר אישור הדוח.
